% ====================================================================
% CAPÍTULO 1: INTRODUCCIÓN Y PLANTEAMIENTO GEOLÓGICO
% ====================================================================
\chapter{Introducción}
% Nota personal: 

El océano desempeña un papel fundamental en la regulación del ciclo global del carbono, participando en el intercambio de dióxido de carbono ($CO_2$) entre la atmósfera y los sistemas marinos \citep{HainSigman2014}. Dentro de este sistema, una de las vías mediante las cuales el carbono puede ser transferido desde las aguas superficiales hacia el océano profundo corresponde a la denominada bomba biológica \citep{HainSigman2014}. Este proceso se encuentra asociado principalmente a la producción de materia orgánica mediante la actividad fotosintética de organismos marinos y a la posterior exportación de una fracción de dicha materia hacia mayores profundidades, donde parte del carbono puede permanecer aislado de la atmósfera durante períodos prolongados \citep{HainSigman2014}.

La eficiencia de la bomba biológica está condicionada por diversos factores físicos, químicos y biológicos, entre ellos la disponibilidad de luz, la temperatura, la circulación y estratificación de la columna de agua, la disponibilidad de nutrientes y la actividad de los organismos que consumen y remineralizan la materia orgánica. Entre los nutrientes esenciales para la productividad primaria destaca el hierro (Fe), un micronutriente cuya disponibilidad puede limitar el crecimiento del fitoplancton incluso en regiones oceánicas donde existen concentraciones relativamente elevadas de macronutrientes, como nitrógeno y fósforo \citep{Tagliabue2017}. El hierro participa en procesos metabólicos fundamentales, como el transporte de electrones durante la fotosíntesis y la respiración celular, y forma parte de proteínas y enzimas involucradas en la utilización de energía y nutrientes. Por ello, su disponibilidad puede influir en la productividad primaria y en la cantidad de materia orgánica producida, aunque la proporción que finalmente alcanza el océano profundo también depende de los procesos de hundimiento, consumo y remineralización \citep{Tagliabue2017}.

El hierro que ingresa a los sistemas oceánicos puede presentar diversos orígenes y encontrarse asociado a diferentes procesos de movilización y transporte. Si bien el aporte de partículas minerales mediante transporte atmosférico ha sido ampliamente estudiado, investigaciones recientes han destacado la importancia de considerar la diversidad de fuentes, caracteristicas mineralogicas y procesos que controlan la disponibilidad del hierro. factores como la composicion mineralogica, el tamaño de las particulas, el grado de alteracion y las transformaciones experimentadas durante el transporte pueden influir sobre la solubilidad y, por esto, sobre la fraccion de hierro potencialmente disponible para los organismos marinos \cite{Zan2025}. De esta manera, la identificacion de una fuente potencial de hierro requiere considerar no solamente su contenido total de este elemento, si no tambien las condiciones que pueden favorecer su movilizacion y posterior transporte.

En este contexto, las características geológicas y ambientales de las zonas continentales pueden constituir un componente relevante para comprender el origen y las posibles rutas de transporte de partículas ricas en hierro hacia sistemas acuáticos. La movilización de material particulado puede estar condicionada por factores como la disponibilidad de sedimentos, las condiciones del suelo y la cobertura vegetal, mientras que su transporte puede depender de variables atmosféricas, topográficas e hidrológicas \citep{Zan2025}. Por lo tanto, la identificación de potenciales fuentes de hierro requiere integrar información proveniente de distintas disciplinas, considerando tanto las características del material como las condiciones ambientales que podrían favorecer su movilización y transporte.

A partir de lo anterior, el presente trabajo busca evaluar el potencial de determinadas fuentes geológicas de hierro y sus posibles rutas de transporte dentro de un área de estudio acotada de la Región del Maule. Para ello, se plantea la integración de antecedentes geológicos, geoquímicos, geomorfológicos, topográficos, hidrológicos y climáticos mediante herramientas de análisis espacial y computacional. El enfoque considera una perspectiva de fuente--transporte--receptor, orientada a identificar espacialmente zonas que presenten condiciones potencialmente favorables para actuar como fuentes de hierro y analizar su posible conectividad con sistemas fluviales y ambientes costeros.

La integración de estas variables permitirá desarrollar una primera aproximación espacial a la distribución de potenciales fuentes de hierro y a los procesos que podrían favorecer su movilización y transporte. Debido a la naturaleza exploratoria del estudio, los resultados serán interpretados en términos de potencialidad y escenarios de movilización, considerando las incertidumbres asociadas a la disponibilidad y calidad de los datos utilizados.

%\section{Planteamiento del Problema}
% Nota: Aquí se explico la contradicción geológica que da origen a la tesis.
%El problema geológico radica en caracterizar el mecanismo... \coment{Esta sección no es necesaria, el planteamiento del problema ya se encuentra en la introducción.}

\section{Hipótesis de Trabajo}%Con formato de hipotesis tal cual H0: hipotesis nula.

\textbf{H$_0$: Hipótesis nula.} La integración de antecedentes geológicos, mineralógicos, geoquímicos, geomorfológicos, topográficos, hidrológicos y climáticos de la Región del Maule no permite distinguir sectores con condiciones relativamente más favorables para constituir fuentes potenciales de hierro ni establecer diferencias espaciales fundamentadas en su posible conectividad con ambientes costeros.

\medskip

\textbf{H$_1$: Hipótesis alternativa.} La integración de antecedentes geológicos, mineralógicos, geoquímicos, geomorfológicos, topográficos, hidrológicos y climáticos de la Región del Maule permite distinguir sectores con condiciones relativamente más favorables para constituir fuentes potenciales de hierro y reconocer posibles conexiones de transporte hacia ambientes costeros.

\section{Objetivos de la Investigación}

\subsection{Objetivo General}
Evaluar exploratoriamente el potencial de fuentes geológicas de hierro y sus posibles rutas de movilización y transporte hacia ambientes costeros de la Región del Maule, mediante la integración de antecedentes geológicos, geoquímicos, geomorfológicos, topográficos, hidrológicos y climáticos.

\subsection{Objetivos Específicos}
% Nota para el formato: Usar la lista con viñetas para ordenar las metas del proyecto.
\begin{itemize}
    \item Caracterizar las principales unidades geológicas y materiales superficiales del área de estudio con potencial para actuar como fuentes de hierro, considerando antecedentes litológicos, mineralógicos, geoquímicos, granulométricos y de alteración.
    \item Construir una base de datos geoespacial que integre antecedentes geológicos, geoquímicos, topográficos, hidrológicos y climáticos relevantes para la identificación y evaluación de fuentes potenciales de hierro.
    \item Analizar la conectividad potencial entre las fuentes identificadas y los posibles receptores mediante rutas de transporte fluvial y atmosférico, utilizando herramientas SIG y computacionales.
\end{itemize}